\section*{第 \ref{ch:g8:ice-on-earth} 章——地球上的冰}

\begin{solution}{exo:g8:ice-on-earth:1}
(a) 积累区。(b) 消融区。(c) 冰舌前端。(d) 一条
终碛。
\end{solution}

\begin{solution}{exo:g8:ice-on-earth:2}
(a) 错：\omterm{def:g8:ice-on-earth:glacier}{冰川}由落在陆地上的雪转化成的冰构成；
冻结的海水是\omterm{def:g8:ice-on-earth:sea-ice}{海冰}。(b) 对。(c) 错：冰仍沿山坡
向下流动，冰舌后退是因为冰融化的速度快于
其流动的速度。(d) 对。
\end{solution}

\begin{solution}{exo:g8:ice-on-earth:3}
(a) \omterm{def:g8:ice-on-earth:glacial-landforms}{U 形谷}。(b) 漂砾。(c) \omterm{def:g8:ice-on-earth:ice-sheet}{冰架}。(d) \omterm{def:g8:ice-on-earth:permafrost}{永久冻土}。
\end{solution}

\begin{solution}{exo:g8:ice-on-earth:4}
(a) \omterm{def:g3:continents-and-oceans:poles}{南极}洲和格陵兰。(b) \omterm{def:g3:continents-and-oceans:poles}{南极}洲，\omterm{def:g4:world-physical-map:altitude}{海平面}约上升
\qty{58}{m}（格陵兰约 \qty{7.4}{m}）。(c) 约十分之一（10~\%）。
\end{solution}

\begin{solution}{exo:g8:ice-on-earth:5}
(a) 1979 年约 705 万~\unit{km^2}，2012 年约 357 万~\unit{km^2}。
(b) $7.05 - 3.57 = 3.48$：约 350 万~\unit{km^2}。
(c) $3.48 \div 7.05 \approx 0.49$：约一半。(d) 某一年的数值取决于
当年夏季的\omterm{def:g1:watching-the-weather:weather}{天气}：2012 年远低于趋势线。
趋势线用到了所有年份，因此某个异常的夏季对它影响很小。
\end{solution}

\begin{solution}{exo:g8:ice-on-earth:6}
(a) $2004 - 1941 = 63$ 年。(b) $12000 \div 63 \approx 190.5$：约
\qty{190}{m} 每年。(c) $2021 - 1984 = 37$ 年；$3200 \div 37 \approx
86$：约 \qty{86}{m} 每年。(d) 缪尔\omterm{def:g8:ice-on-earth:glacier}{冰川}，速度快一倍以上。
\end{solution}

\begin{solution}{exo:g8:ice-on-earth:7}
在\emph{积累区}，即山地高处，降雪量
每年都多于融化量。新雪覆盖并压实旧雪，旧雪
约一年后变为\emph{粒雪}，随后在更深处并经过更长时间，
变为坚冰。冰在自身重量下，缓慢\emph{流动}，
沿\omterm{def:g1:landscapes-around-us:landform}{山谷}向下。在下游的\emph{消融区}，冰的融化
速度快于降雪补充的速度；到达冰舌前端的冰，
是在高处形成的冰，往往形成于很久以前。
\end{solution}

\begin{solution}{exo:g8:ice-on-earth:8}
(a) 烧杯 B 中：融化的冰流入水中并使水位升高；在
烧杯 A 中水位保持不变。(b) 烧杯 A 代表\omterm{def:g8:ice-on-earth:ice-sheet}{冰山}（已漂浮
在海中的冰）；烧杯 B 代表\omterm{def:g8:ice-on-earth:ice-sheet}{冰盖}（陆地上的冰）。
(c) \omterm{def:g8:ice-on-earth:ice-sheet}{冰架}挡住了其后方的\omterm{def:g8:ice-on-earth:glacier}{冰川}。\omterm{def:g8:ice-on-earth:ice-sheet}{冰架}破裂时，
这些\omterm{def:g8:ice-on-earth:glacier}{冰川}加速，将更多陆地冰注入\omterm{def:g2:water-on-earth:ocean}{海洋}，
从而使\omterm{def:g4:world-physical-map:altitude}{海平面}升高。
\end{solution}

\begin{solution}{exo:g8:ice-on-earth:9}
(a) \omterm{def:g8:ice-on-earth:glacial-landforms}{U 形谷}，由\omterm{def:g8:ice-on-earth:glacier}{冰川}\omterm{def:g5:water-shapes-the-land:erosion}{侵蚀}形成。(b) \omterm{def:g8:ice-on-earth:glacier}{冰川}将其从
\omterm{def:g1:landscapes-around-us:landform}{山谷}上游（那里有花岗岩）搬运而来，冰融化时将其留下，
因此它是一块漂砾。(c) 侧碛：其高度表明\omterm{def:g8:ice-on-earth:glacier}{冰川}
曾经有多厚；冰至少曾达到谷壁上的
这一高度。
\end{solution}

\begin{solution}{exo:g8:ice-on-earth:10}
(a) $7.4 \div 10 = \qty{0.74}{m}$。(b) $58 \div 100 = \qty{0.58}{m}$。
(c) $58 \div 7.4 \approx 7.8$：约多 8 倍。
\end{solution}

\begin{solution}{exo:g8:ice-on-earth:11}
(a) $7.21 \times 10^6 = \qty{7210000}{km^2}$，$4.65 \times 10^6 =
\qty{4650000}{km^2}$。(b) $7.21 - 4.65 = \num{2.56}\text{ 百万 }\unit{km^2}$，即
\qty{2560000}{km^2}。(c) $2.56 \div 3 \approx 0.85$：约
每十年 85 万~\unit{km^2}，略多于趋势线的
76 万~\unit{km^2}；每十年的平均值仍包含该十年
自身十个夏季的起伏，但两种读图方式都显示出相同的
巨大损失。(d) 范围不可能降至零以下，而且趋势线
只描述已测量的年份：今后的情况取决于\omterm{def:g3:continents-and-oceans:poles}{北极}
\omterm{def:g5:weather-and-climate:climate}{气候}如何变化，因此该直线并非预测。
\end{solution}

\begin{solution}{exo:g8:ice-on-earth:12}
(a) $120 \div 1 = 120$：当时\omterm{def:g1:landscapes-around-us:water}{海岸}比现在
向海外推进约 \qty{120}{km}。(b) 水汽从\omterm{def:g2:water-on-earth:ocean}{海洋}蒸发，以雪的形式落在
\omterm{def:g3:continents-and-oceans:continent}{大陆}上，并以大\omterm{def:g8:ice-on-earth:ice-sheet}{冰盖}冰的形式留在那里。(c) 终碛
形成于\omterm{def:g8:ice-on-earth:glacier}{冰川}末端。远离任何\omterm{def:g4:world-physical-map:range}{山脉}的\omterm{def:g1:landscapes-around-us:landform}{平原}上的一道终碛，
说明\omterm{def:g8:ice-on-earth:ice-sheet}{冰盖}曾抵达该\omterm{def:g1:landscapes-around-us:landform}{平原}，当时
冰覆盖的陆地比今天多得多。
\end{solution}

\begin{solution}{pb:g8:ice-on-earth:1}
\textbf{1.} $1080 - 0 = \qty{1080}{m}$。

\textbf{2.} 1960--1970 年：\qty{180}{m}；1970--1980 年：\qty{150}{m}；
1980--1990 年：\qty{190}{m}；1990--2000 年：\qty{180}{m}；2000--2010 年：
\qty{180}{m}；2010--2020 年：\qty{200}{m}。最少的是 1970--1980 年，最多的是
2010--2020 年。

\textbf{3.} 否：后退中的\omterm{def:g8:ice-on-earth:glacier}{冰川}仍沿山坡向下流动，其冰舌
后退是因为冰融化的速度快于冰流来的速度。

\textbf{4.} 终碛：冰将岩屑沿\omterm{def:g1:landscapes-around-us:landform}{山谷}向下搬运，
并在冰舌前端融化处将其堆积下来，
此时冰舌仍停留在该处。

\textbf{5.} $1080 \div 6 = \qty{180}{m}$，即每十年。

\textbf{6.} $180 \div 10 = \qty{18}{m}$，即每年。

\textbf{7.} $190 \div 18 \approx 10.6$：约快 11 倍。

\textbf{8.} 较少：其消融区损失的冰多于积累区获得的冰，
因此冰舌后退。

\textbf{9.} U 形：谷底宽阔、谷壁陡峭，因为\omterm{def:g8:ice-on-earth:glacier}{冰川}
在充填\omterm{def:g1:landscapes-around-us:landform}{山谷}的同时拓宽、加深了\omterm{def:g1:landscapes-around-us:landform}{山谷}。

\textbf{10.} $1980 - 1080 = \qty{900}{m}$。

\textbf{11.} $900 \div 18 = 50$ 年，即在 \textbf{2070} 年。

\textbf{12.} 速度逐十年变化（从 150
到 \qty{200}{m}），未来的后退取决于未来几年的降雪量与夏季气温：
过去的速率并非预测。
\end{solution}
